\documentclass[../../main.tex]{subfiles}
% \documentclass[a4paper,10pt]{ltjsarticle}
\begin{document}
\rhead{京都大学後期$2005$年$6$番}
\setcounter{equation}{0}
\begin{oframed}
  \subfile{problem}
\end{oframed}
\setlength{\columnseprule}{0.4pt}
\begin{multicols}{2}
  {\bf [解答]}
  表の出た100円玉の枚数$X$, 500円玉の枚数$Y$とすると，
  $X \le Y-1$ となる確率$P(X \le Y-1)$を求めれば良い.
  表と裏の対称性から，「表の出た $100$ 円玉の枚数より表の出た $500$ 円玉の枚
  数の方が多い」確率は，表裏を入れ替えた
  「裏の出た $100$ 円玉の枚数より裏の出た $500$ 円玉の枚数の方が多い」
  確率と同じである．裏の出た枚数は100円玉$n-X$枚，500円玉は$n+1-Y$だから
  \begin{align}
    P(X \le Y-1)
    &= P(n-X \le (n+1-Y)-1) \\
    &= P(n-X \le n-Y) \\
    &= P(Y \le X) \label{eq:1}
  \end{align}
  である．

  \begin{figure}[H]
    \centering
    \label{2005-6:fig:1}
    \begin{tikzpicture}
      \draw[->] (0,0) -- (7,0) node[right] {$x$};
      \draw[->] (0,0) -- (0,7) node[above] {$y$};

      % 目盛
      \foreach \x in {1,2,3,4,5} \draw (\x,0) -- (\x,0.1);
      \foreach \y in {1,2,3,4,5} \draw (0,\y) -- (0.1,\y);

      \node at (5,-0.3) {$n$};
      \node at (-0.3,6) {$n+1$};
      \node at (-0.3,1) {$1$};

      \draw (0,0) rectangle (5,6);

      \draw[thick, dashed] (0,0) -- (5,5) node[below right] {$Y=X$};
      \draw[thick, dashed] (0,1) -- (5,6) node[above left] {$Y=X+1$};

      % Y >= X+1 の領域（上側）
      \fill[pattern=north east lines, pattern color=blue!70!black] (0,1) -- (0,6) -- (5,6) -- cycle;
      % X >= Y の領域（下側）
      \fill[pattern=north west lines, pattern color=gray!70!black] (0,0) -- (5,0) -- (5,5) -- cycle;
    \end{tikzpicture}
    \caption{$X$と$Y$のとりうるパターンを平面上に図示した様子．}
  \end{figure}

  ここで，$P(X \le Y-1)$と$P(Y \le X)$の二つのパターンは枚数に関して全てのパターンを尽くしており，
  \begin{align}
    P(X \le Y-1) + P(Y \le X) = 1 \label{eq:2}
  \end{align}
  である．この様子を\cref{2005-6:fig:1}に示す．

  \cref{eq:1,eq:2}より
  \begin{align}
    P(X+1 \le Y) = \dfrac{1}{2}
  \end{align}
  が求める答えである．

  {\bf [別解1]}
  1玉の500円玉に注目し, その表裏で場合分けする.
  残り$n$枚ずつの$100$円玉，$500$円玉の表の枚数をそれぞれ $X, Y$ とする.

  \textbf{$1^\circ$ 表の時}
  残り$n$枚ずつ投げて, $X \le Y$ となれば良く, 確率 $P(X \le Y)$

  \textbf{$2^\circ$ 裏の時}
  残り$n$枚ずつ投げて, $X < Y$ となれば良く, 確率 $P(X < Y)$

  $1$枚の$500$円玉を除くとコインが$n$枚ずつあるため対称性から $P(X < Y) = P(X > Y)$ だから求める確率は
  \begin{align}
    \dfrac{1}{2} \left( P(X \le Y) + P(X < Y) \right) = \dfrac{1}{2} \left( P(X \le Y) + P(X > Y) \right)
  \end{align}
  であり，右辺$P(X \le Y) + P(X > Y)$は全ての事象を尽くすから$1$なので求める確率は
  \begin{align}
    \dfrac{1}{2} \left( P(X \le Y) + P(X < Y) \right) = \dfrac{1}{2}
  \end{align}
  である．

  {\bf [別解2]}
  500円が$k$枚表である確率を$P_k$とおくと,
  100円が$j$枚表である確率を$q_j$とおくと,
  \begin{align}
    P_k & = {}_{n+1}\mathrm{C}_k \left(\dfrac{1}{2}\right)^{n+1} \\
    q_j & = {}_{n}\mathrm{C}_j \left(\dfrac{1}{2}\right)^{n}
  \end{align}
  で，
  \begin{align}
    P & = \sum_{k=1}^{n+1} P_k (q_0 + \dots + q_{k-1})                                                                                                 \\
    & = \left(\dfrac{1}{2}\right)^{2n+1} \sum_{k=1}^{n+1} {}_{n+1}\mathrm{C}_k ({}_n\mathrm{C}_0 + \dots + {}_n\mathrm{C}_{k-1}) \label{eq:3}
  \end{align}
  ここで, $A = \sum_{k=1}^{n+1} {}_{n+1}\mathrm{C}_k ({}_n\mathrm{C}_0 + \dots + {}_n\mathrm{C}_{k-1})$ とおくと, $A$ は $(1+x)^n (1+x)^{n+1}$ の展開項のうち, $x^{n+1}$ 次以上のもの和である.
  \begin{align}
    (x+1)^{2n+1} = {}_{2n+1}\mathrm{C}_0 x^0 + \dots + {}_{2n+1}\mathrm{C}_n x^n + {}_{2n+1}\mathrm{C}_{n+1} x^{n+1} + \dots + {}_{2n+1}\mathrm{C}_{2n+1} x^{2n+1}
  \end{align}
  として, ${}_{2n+1}\mathrm{C}_{2n+1-k} = {}_{2n+1}\mathrm{C}_k$ だから.
  \begin{align}
    2^{2n+1} = 2 \left( {}_{2n+1}\mathrm{C}_0 + \dots + {}_{2n+1}\mathrm{C}_n \right) = 2A
  \end{align}
  \begin{align}
    \therefore A = 2^{2n} \label{eq:4}
  \end{align}
  \cref{eq:3,eq:4}から,
  \begin{align}
    P = \left(\dfrac{1}{2}\right)^{2n+1} 2^{2n} = \dfrac{1}{2} \cdots\text{(答)}
  \end{align}
  を得る．

  {\bf [解説]}

  \newpage
\end{multicols}
\end{document}
